\uuid{7M38}
\titre{Équation homogène/solution particulière}
\auteur{hanani ; mzari}
\organisation{exo7}
\licence{CC BY-NC-SA 4.0}
\datecreate{2020-01-20}
\niveau{L1}
\module{Analyse}
\chapitre{Equation différentielle}
\sousChapitre{Résolution d'équation différentielle du premier ordre}
\difficulte{2}
\theme{Equations du premier ordre, équation homogène, solution particulière}
\exo7qcmid{lille-2.equations-differentielles.006}

\contenu{
\begin{question}
\qclassification{225.01}
\qtags{motcle=équation homogène/solution particulière}
On considère l'équation différentielle $(E)$ : $\displaystyle x^2y'-(2x-1)y=x^2$ sur $\Rr$. Parmi les affirmations suivantes cocher celles qui sont vraies :
\begin{answers}
\good{La fonction $\displaystyle y=x^2$ est une solution de $(E)$ sur $\Rr$.}
\bad{La fonction $\displaystyle y=x^2\left(1-\mathrm{e}^{1/x}\right)$ est une solution de l'équation homogène.}
\bad{La fonction $\displaystyle y=2x^2\mathrm{e}^{1/x}$ est une solution de $(E)$ sur $]0,+\infty[$.}
\good{La fonction $\displaystyle y=x^2\left(1-\mathrm{e}^{1/x}\right)$ est une solution de $(E)$ sur $]0,+\infty[$.}
\end{answers}
\begin{explanations}
La fonction $y_0=x^2$ est continue, dérivable sur $\Rr$ et elle vérifie $(E)$ ; c'est une solution de $(E)$ sur $\Rr$. De même, la fonction $\displaystyle y=x^2\left(1-\mathrm{e}^{1/x}\right)$ est continue, dérivable sur $]0,+\infty[$; c'est une solution de $(E)$ sur $]0,+\infty[$.
\end{explanations}
\end{question}
}
